\documentclass[conference,letterpaper]{IEEEtran}
\usepackage[T1]{fontenc}
\usepackage{amsmath,amssymb}
\usepackage{newtxtext,newtxmath}
\usepackage{microtype}
\usepackage{booktabs}
\usepackage{graphicx}
\usepackage{cite}
\usepackage{placeins}
\usepackage{needspace}
\usepackage{balance}
\usepackage{xcolor}
\usepackage[hidelinks]{hyperref}
\hypersetup{pdftitle={Placement-Aware Refinement for CLIP-Based Weakly Supervised Video Anomaly Detection},pdfauthor={Anonymous Authors},pdfsubject={Weakly supervised video anomaly detection},pdfkeywords={WSVAD, CLIP, temporal variation, semantic adaptation}}
\newcommand{\Agg}{\operatorname{Agg}}
\newcommand{\Norm}{\operatorname{norm}}
\newcommand{\softmax}{\operatorname{softmax}}
\newcommand{\BCE}{\operatorname{BCE}}
\newcommand{\mlp}{\mathrm{mlp}}
\hyphenation{VadCLIP anomalous anomaly aggregation}
\begin{document}
\title{Placement-Aware Refinement for CLIP-Based Weakly Supervised Video Anomaly Detection}
\author{\IEEEauthorblockN{Anonymous Authors}}
\maketitle
\begin{abstract}
Fixed-length temporal aggregation in CLIP-based weakly supervised video anomaly detection preserves average appearance but suppresses local feature variation. Additional semantic supervision can also change classification predictions through their shared representation. We propose a placement-aware refinement framework that addresses these two issues by controlling where temporal evidence enters the model and where semantic optimization acts. First, pre-aggregation variation refinement extracts differences from raw CLIP snippet features and injects aligned variation statistics through a zero-initialized low-rank residual. A sparse-neighborhood refinement further strengthens similarity-based temporal propagation. Second, pathway-restricted semantic adaptation optimizes a semantic multiple-instance objective through the existing text-side adapter, preserving classification outputs by construction. Under fixed-epoch evaluation on UCF-Crime and XD-Violence, the method improves the matched VadCLIP baseline by 0.41 AUC and 1.66 AP percentage points, respectively. Under the released VadCLIP training settings, the three-seed paired gains are $0.35\pm0.03$ AUC and $0.73\pm0.16$ AP percentage points on UCF-Crime and XD-Violence, respectively, with best-run results of 88.25\% AUC and 84.94\% AP. Component and placement ablations validate both refinements. The method adds 132,736 parameters and increases neural scoring latency by 3.3--3.6\%, with no additional inference operations from semantic adaptation.
\end{abstract}
\begin{IEEEkeywords}
weakly supervised video anomaly detection, vision-language models, temporal variation, semantic adaptation, intelligent surveillance
\end{IEEEkeywords}
% Introduction. Numbers via macros; claims follow evidence/discrepancies.md.

\section{Introduction}
\label{sec:intro}

Weakly supervised video anomaly detection learns to localize abnormal events in time from video-level labels alone. Current methods start from frozen snippet features, most often from CLIP~\cite{radford2021clip}, and train a light temporal model on top with a multiple-instance objective~\cite{sultani2018real}. VadCLIP~\cite{wu2024vadclip} extends this design with a second output that aligns snippets with learnable class prompts, and it is a common base for recent work~\cite{yang2024tpwng,wu2024stprompt}.

Two design choices in this pipeline limit what the model can learn. The first concerns the input. Before any learning, each training video is resampled to a fixed length, and long videos are mean-pooled over uniform segments. On the training sets we use, this affects \ucfTrainAvgFrac\% of UCF-Crime and \xdTrainAvgFrac\% of XD-Violence items. Changes between snippets inside a segment do not survive pooling, and differencing the pooled sequence only measures change between segment means. The second concerns the objective. VadCLIP's semantic output is computed downstream of its classification output, and both read the same representation. Adding a semantic multiple-instance loss to the whole network therefore also moves the classification output. In a single-seed run, doing so raised XD-Violence AP by \fullGXdApTwo{} points and changed UCF-Crime AUC by \fullGUcfAucOne{} points.

We address the two choices in two training stages. Stage~1 computes the absolute difference between consecutive raw snippets before pooling, aggregates it with the same segment boundaries as the appearance features, and injects it through a zero-initialized low-rank residual; it also mixes VadCLIP's similarity graph with a sparse $k$-nearest-neighbour graph. Stage~2 optimizes the semantic multiple-instance loss with respect to VadCLIP's text-side adapter only. That adapter lies downstream of the classification output, which is therefore unchanged by construction.

We evaluate under a protocol in which neither training nor checkpoint selection uses the test set, retrain VadCLIP under the same protocol, and report paired differences over three seeds. Our method improves UCF-Crime AUC by \ucfGainMean{} points and XD-Violence AP by \xdGainMean{} points, with the same sign on every seed. Computing the same variation after pooling instead falls below the baseline on every seed, and an equal-budget control without the semantic objective shows that the Stage~2 gain is not an effect of additional training. Attribution experiments locate the UCF-Crime gain in Stage~1 and the XD-Violence gain in Stage~2.

Our contributions are as follows.
\begin{itemize}
  \item We show that computing temporal variation before fixed-length pooling, rather than after it, changes what a CLIP-based detector learns from long training videos, and we inject it through a residual that adds \pvrParams{} parameters.
  \item We restrict semantic post-training to the part of VadCLIP downstream of its classification output, which keeps that output unchanged while improving the semantic output, without adding parameters.
  \item We evaluate both stages in a matched multi-seed comparison with an equal-budget control and component attribution, under a protocol in which neither training nor checkpoint selection uses the test set.
\end{itemize}

% Related work. Only works listed in refs/citation_ledger.csv are cited.

\section{Related Work}
\label{sec:related}

\textbf{WSVAD and CLIP-based detectors.} Most WSVAD methods treat a video as a bag of snippets and train a scorer with a multiple-instance objective over the highest-scoring snippets~\cite{sultani2018real}. Later work refines the snippet representation through feature-magnitude learning~\cite{tian2021rtfm}, memory units with uncertainty regulation~\cite{zhou2023urdmu}, and similarity and distance graphs for temporal context~\cite{wu2020xdviolence}. With CLIP~\cite{radford2021clip} features, CLIP-TSA adds temporal self-attention~\cite{joo2023cliptsa}, and VadCLIP~\cite{wu2024vadclip} adds a semantic output that matches snippets against learnable class prompts. Follow-up work enriches the prompts with normality guidance~\cite{yang2024tpwng} or spatio-temporal cues~\cite{wu2024stprompt}. These methods all model time after the fixed-length resampling of snippet features. We change what reaches that step and which parameters a semantic objective may update.

\textbf{Temporal modelling and lightweight adaptation.} Explicit temporal differences are an established cue in action recognition, where TDN computes frame differences inside a trainable backbone~\cite{wang2021tdn} and TSM shifts channels across time~\cite{lin2019tsm}. In WSVAD the backbone is frozen, and the only temporal operation before learning is aggregation; our refinement targets that step. Parameter-efficient methods insert small modules or low-rank updates into a frozen model~\cite{houlsby2019adapter,hu2021lora}, and CLIP-Adapter~\cite{gao2024clipadapter} applies the idea to CLIP. Instead of adding a module, our second design restricts an extra objective to an existing adapter selected by its position in the prediction pathways.

% Method. Symbols follow evidence/formula_code_audit.md.

\section{Method}
\label{sec:method}

We build on VadCLIP~\cite{wu2024vadclip}, which takes frozen CLIP snippet features $X \in \mathbb{R}^{L \times d}$ ($d=512$) and maps each training video to $T=256$ steps with an operator $\mathrm{Agg}$: longer videos are mean-pooled over $T$ uniform segments and shorter ones are zero-padded. A temporal transformer and two graph branches produce a shared representation $h$. The classification output $s^{(1)} = f_{\mathrm{cls}}(h + \mathrm{mlp}_2(h))$ gives snippet anomaly logits. The semantic output pools $h$ with weights from $s^{(1)}$ into a video vector $a$, refines each class prompt with a text-side adapter, $t_c = (t^0_c + a) + \mathrm{mlp}_1(t^0_c + a)$, and scores $s^{(2)}_{j,c} = \cos(h_j, t_c)/\tau$, from which the anomaly probability $p_j = 1 - \mathrm{softmax}(s^{(2)}_j)_{\mathrm{normal}}$ follows. Following VadCLIP, UCF-Crime is reported with $s^{(1)}$ and XD-Violence with $p$. Since $s^{(1)}$ does not depend on $\mathrm{mlp}_1$, the semantic pathway can be adapted without touching the classification pathway.

Our method applies refinement at two positions (Fig.~\ref{fig:overview}): temporal variation is extracted before aggregation (Stage~1), and semantic adaptation is applied after the pathways separate (Stage~2).

% F1 (6-page version): two panels. Blue = added in Stage 1, orange = updated in Stage 2, grey = VadCLIP (frozen in Stage 2).
\begin{figure*}[t]
\centering
\begin{tikzpicture}[
  font=\footnotesize, >=Stealth,
  box/.style={draw, rounded corners=2pt, minimum height=6.5mm, inner sep=2.5pt, align=center},
  new/.style={box, fill=blue!12},
  trn/.style={box, fill=orange!35, very thick},
  frz/.style={box, fill=gray!12},
  op/.style={draw, circle, inner sep=0.8pt, fill=white},
  tag/.style={font=\footnotesize\bfseries, inner sep=1.5pt},
  lab/.style={font=\scriptsize, inner sep=1pt}]

  % ---------------- Panel (a) ----------------
  \node[frz] (X) at (0,0) {Raw\\snippets $X$};
  \node[new] (D) at (2.1,0.75) {$\lvert X_t - X_{t-1}\rvert$};
  \node[frz] (AggA) at (3.9,-0.75) {$\mathrm{Agg}$};
  \node[new] (AggV) at (3.9,0.75) {$\mathrm{Agg}$};
  \node[new] (R) at (5.5,0.75) {zero-init\\residual};
  \node[op] (plus) at (5.5,-0.75) {$+$};
  \node[frz] (Enc) at (7.3,-0.75) {temporal\\model};
  \node[box] (h) at (8.6,-0.75) {$h$};
  \draw[->] (X.east) -- ++(0.25,0) |- (D.west);
  \draw[->] (X.east) -- ++(0.25,0) |- (AggA.west);
  \draw[->] (D) -- (AggV);
  \draw[->] (AggV) -- node[lab, above] {$V$} (R);
  \draw[->] (R) -- node[lab, right] {$\beta$} (plus);
  \draw[->] (AggA) -- node[lab, above] {$A$} (plus);
  \draw[->] (plus) -- node[lab, above] {$Z$} (Enc);
  \draw[->] (Enc) -- (h);
  % sparse kNN callout
  \node[new, dashed, font=\scriptsize] (knn) at (7.3,0.65) {sparse $k$NN\\graph mix};
  \draw[->, dashed] (knn) -- (Enc);
  % emphasis: before aggregation
  \begin{scope}[on background layer]
    \node[draw=blue!60, thick, dashed, rounded corners=3pt, fit=(D)(AggV), inner sep=3pt,
          label={[tag, text=blue!60!black]above:Before Aggregation}] {};
  \end{scope}
  \node[font=\small] at (4.3,-2.15) {(a) Earlier temporal refinement (Stage~1)};

  % ---------------- Panel (b) ----------------
  \begin{scope}[xshift=10.3cm]
    \node[box] (h2) at (0,-0.1) {shared\\$h$};
    \coordinate (split) at (0.9,-0.1);
    \node[frz] (cls) at (2.6,0.75) {classification branch\\$f_{\mathrm{cls}}, \mathrm{mlp}_2$ [frozen]};
    \node[box] (s1) at (4.9,0.75) {$s^{(1)}$};
    \node[trn] (mlp1) at (2.6,-0.95) {semantic branch\\$\mathrm{mlp}_1$ [updated]};
    \node[box] (p) at (4.9,-0.95) {$s^{(2)}\!\rightarrow p$};
    \draw[-] (h2) -- (split);
    \draw[->] (split) |- (cls.west);
    \draw[->] (split) |- (mlp1.west);
    \draw[->] (cls) -- (s1);
    \draw[->] (mlp1) -- (p);
    \draw[->, gray] (s1.south) -- node[lab, right, text=gray!70!black] {pooling $a$} (s1.south |- mlp1.north east);
    \node[lab, text=orange!70!black, align=left] (loss) at (6.6,-0.95) {$\mathcal{L}_{\mathrm{CLASM}}$\\$+\lambda\mathcal{L}_{G}$};
    \draw[->, dashed, orange!70!black] (p) -- (loss);
    \draw[->, dashed, orange!70!black] (loss.south) -- ++(0,-0.3) -| node[lab, pos=0.25, below] {gradient to $\mathrm{mlp}_1$ only} (mlp1.south);
    \node[tag, text=orange!60!black, fill=white, draw=orange!60, dashed, rounded corners=2pt] at (0.9,1.35) {After Pathway Separation};
    \draw[orange!60, dashed] (0.9,1.1) -- (split);
    \node[font=\small] at (3.2,-2.15) {(b) Later semantic adaptation (Stage~2)};
  \end{scope}
\end{tikzpicture}
\caption{Overview. (a) Stage~1 extracts snippet differences \emph{before} fixed-length aggregation, pools them with the same segment boundaries as the appearance features, and fuses them through a zero-initialized residual; a sparse $k$NN graph mix refines the similarity branch. (b) Stage~2 restricts the semantic objective to $\mathrm{mlp}_1$, which lies \emph{after} the classification and semantic pathways separate, so $s^{(1)}$ is unchanged. Blue: added in Stage~1; orange: updated in Stage~2; grey: VadCLIP, frozen in Stage~2.}
\label{fig:overview}
\end{figure*}


\subsection{Pre-Aggregation Temporal Refinement}
\label{sec:pvr}

For a long training video, $\mathrm{Agg}$ averages on average \ucfBinSize{} (UCF-Crime) and \xdBinSize{} (XD-Violence) snippets per segment, so differences of the aggregated sequence only measure change between segment means. We instead compute variation on the raw snippets,
\begin{equation}
  D_1 = 0, \qquad D_t = \lvert X_t - X_{t-1} \rvert, \quad t = 2, \dots, L,
  \label{eq:diff}
\end{equation}
and aggregate it with the same segment boundaries as appearance, $A = \mathrm{Agg}(X)$ and $V = \mathrm{Agg}(D)$, so both streams stay aligned at training and test time. The variation enters through a low-rank residual,
\begin{equation}
  Z = A + \beta\, W_2\, \mathrm{GELU}\big(W_1\, \mathrm{LN}(V)\big),
  \label{eq:pvr}
\end{equation}
with $W_1 \in \mathbb{R}^{128 \times d}$, $W_2 \in \mathbb{R}^{d \times 128}$, and $\beta = 0.2$. $W_2$ and its bias are zero-initialized, so training starts exactly from VadCLIP. $Z$ replaces $A$ as the input of the temporal model and therefore reaches both outputs. This design preserves within-segment variation at a cost of \pvrParams{} parameters.

\textbf{Sparse neighbourhood refinement.} As a supporting refinement within Stage~1, we sharpen VadCLIP's similarity graph. Its thresholded softmax still assigns weight $e^0$ to every zeroed entry, so each snippet stays connected to all others. We mix it with a sparse $k$-nearest-neighbour graph, $G = (1-\alpha)\,G_{\mathrm{orig}} + \alpha\,\mathrm{softmax}(M_k(S))$, where $S$ is the cosine similarity matrix and $M_k$ keeps the $k$ largest entries per row; we use $k=16$ and $\alpha=0.25$. No parameters are added.

\subsection{Pathway-Restricted Semantic Adaptation}
\label{sec:isa}

VadCLIP trains $s^{(2)}$ with the class-level CLASM objective. We add a binary multiple-instance target on $p$: for a video of valid length $\ell$ and label $y$, let $q$ be the mean of the $\lfloor \ell/16 \rfloor + 1$ largest $p_j$, and
\begin{equation}
  \mathcal{L} = \mathcal{L}_{\mathrm{CLASM}} + \lambda\, \mathrm{BCE}(q, y), \qquad \lambda = 0.1.
  \label{eq:isa}
\end{equation}
Optimized over the whole network, this objective also changes $h$ and hence $s^{(1)}$. Stage~2 starts from the Stage~1 model and optimizes Eq.~\eqref{eq:isa} with respect to $\mathrm{mlp}_1$ only, with all other parameters frozen. Because $s^{(1)}$ is independent of $\mathrm{mlp}_1$, the restricted update prevents any change to the classification pathway while the semantic output adapts. Stage~2 adds no parameters and no inference cost.

\subsection{Implementation Details}
\label{sec:impl}

Stage~1 follows VadCLIP's released configuration for each dataset and trains for 10 epochs. Stage~2 trains the \mlpOneParams{} parameters of $\mathrm{mlp}_1$ for 3 epochs with AdamW (learning rate $2\times10^{-5}$, weight decay $0.01$), batch size 64 on UCF-Crime and 96 on XD-Violence. The hyperparameters $\beta$, $k$, $\alpha$, and $\lambda$ are shared across datasets.

% Experiments. All numbers come from numbers.tex / numbers_results.tex and tables/*.tex.

\section{Experiments}
\label{sec:experiments}

\subsection{Setup}
\label{sec:setup}

We use UCF-Crime~\cite{sultani2018real} (290 test videos) and XD-Violence~\cite{wu2020xdviolence} (800 test videos) with the CLIP features, splits, and frame-level ground truth released with VadCLIP. Following VadCLIP, we report frame-level AUC of $s^{(1)}$ on UCF-Crime and frame-level AP of $p$ on XD-Violence. Controlled experiments and ablations use a test-independent fixed-epoch protocol with matched seeds: every model is trained for a fixed number of epochs and the final checkpoint is evaluated. Each condition is run with seeds 234, 235, and 236, which share initialization and sample order across arms, and we report mean$\pm$SD, the paired difference per seed, and the number of positive seeds. We additionally evaluate the final method using VadCLIP's official training script.

\subsection{Main Results}
\label{sec:main}

% Auto-generated by scripts/make_tables.py. Do not edit by hand.
\begin{table}[t]
\centering
\small
\setlength{\tabcolsep}{4pt}
\caption{Results on UCF-Crime (AUC, \%) and XD-Violence (AP, \%) with CLIP features. Published rows are taken from VadCLIP~\cite{wu2024vadclip} and serve as benchmark context; each follows its own training protocol. Method claims rest on the matched comparison: mean$\pm$SD over three seeds, paired $\Delta$ per seed, number of positive seeds in parentheses. The bottom block repeats the matched comparison under VadCLIP's official script, which selects the checkpoint on the test set; per-seed values are in Table~\ref{tab:official}. Only rows within a block are directly comparable.}
\label{tab:main}
\begin{tabular}{lcc}
\toprule
Method & UCF AUC & XD AP \\
\midrule
RTFM~\cite{tian2021rtfm} & 85.66 & 78.27 \\
UR-DMU~\cite{zhou2023urdmu} & 86.75 & 82.41 \\
CLIP-TSA~\cite{joo2023cliptsa} & 87.58 & 82.17 \\
VadCLIP~\cite{wu2024vadclip} & 88.02 & 84.51 \\
\midrule
\multicolumn{3}{l}{\emph{Matched comparison (same protocol and seeds)}} \\
VadCLIP, retrained & 86.25$\pm$0.83 & 77.57$\pm$2.19 \\
Ours & 86.53$\pm$0.85 & 79.23$\pm$1.87 \\
Paired $\Delta$ & +0.28$\pm$0.15 (3/3) & +1.66$\pm$0.34 (3/3) \\
\midrule
\multicolumn{3}{l}{\emph{VadCLIP's official script (best-epoch selection)}} \\
VadCLIP, retrained & -- & 82.91$\pm$1.23 \\
Ours & -- & 83.63$\pm$1.18 \\
\bottomrule
\end{tabular}
\end{table}


Table~\ref{tab:main} places our results in the context of published CLIP-feature methods and reports the matched comparison. Against VadCLIP retrained under the same protocol and seeds, our method improves UCF-Crime AUC by \ucfGainMean{}$\pm$\ucfGainSd{} points and XD-Violence AP by \xdGainMean{}$\pm$\xdGainSd{} points, and the paired gain is positive on every seed for both datasets. The two gains come from different designs: Stage~2 leaves $s^{(1)}$ unchanged, so the UCF-Crime gain is produced by Stage~1, whereas Stage~1 alone changes XD-Violence AP by only \xdStageOneMean{} points and the XD-Violence gain is produced by Stage~2.

\subsection{Mechanism Ablations}
\label{sec:ablation}

% Auto-generated by scripts/make_tables.py. Do not edit by hand.
\begin{table}[t]
\centering
\small
\setlength{\tabcolsep}{4pt}
\caption{Mechanism ablations (\%, mean$\pm$SD over three matched seeds). $\Delta$ is the paired difference to VadCLIP (B0) unless stated otherwise. (a) Post-aggregation computes the difference on the pooled sequence with all other settings fixed. (b) C runs the same budget from the same checkpoint with $\lambda=0$.}
\label{tab:ablation}
\begin{tabular}{lcc}
\toprule
Variant & Score & $\Delta$ \\
\midrule
\multicolumn{3}{l}{\emph{(a) Stage 1, UCF-Crime AUC}} \\
VadCLIP (B0) & 86.25$\pm$0.83 & -- \\
Post-aggregation variation & 85.86$\pm$0.84 & $-$0.39$\pm$0.10 (0/3) \\
Pre-aggregation variation & 86.53$\pm$0.85 & +0.28$\pm$0.15 (3/3) \\
\quad Pre $-$ Post &  & +0.67$\pm$0.05 (3/3) \\
\midrule
\multicolumn{3}{l}{\emph{(b) Stage 2, XD-Violence AP}} \\
VadCLIP (B0) & 77.57$\pm$2.19 & -- \\
Stage 1 & 77.44$\pm$2.52 & $-$0.14$\pm$0.33 (1/3) \\
+ CLASM only (C) & 77.59$\pm$2.28 & +0.02$\pm$0.30 (2/3) \\
+ Restricted objective & 79.23$\pm$1.87 & +1.66$\pm$0.34 (3/3) \\
\quad Restricted $-$ C &  & +1.64$\pm$0.43 (3/3) \\
B0 + restricted objective & 79.46$\pm$1.48 & +1.89$\pm$0.76 (3/3) \\
\bottomrule
\end{tabular}
\end{table}


\textbf{Where temporal variation is extracted.} Table~\ref{tab:ablation}(a) computes the same difference after aggregation, $D'_j = \lvert A_j - A_{j-1} \rvert$, with every other setting fixed. This comparison isolates the position of the difference relative to pooling. Post-aggregation variation falls below the baseline on every seed (\postBaseMean{} points), while pre-aggregation variation exceeds it by \prePostMean{}$\pm$\prePostSd{} AUC points (\prePostPos{} seeds). The benefit therefore comes from preserving within-segment change, not from adding a difference stream as such.

\textbf{Where semantic optimization is applied.} Table~\ref{tab:ablation}(b) compares Stage~2 with an equal-budget control C that runs the same epochs from the same checkpoint and sample order with $\lambda=0$. The restricted semantic objective adds \tMinusCMean{}$\pm$\tMinusCSd{} AP points over C (\tMinusCPos{} seeds), so the gain comes from the objective rather than from extra training. Applied directly to the retrained VadCLIP, Stage~2 improves AP by \bzTMean{}$\pm$\bzTSd{} points (\bzTPos{} seeds), which attributes the XD-Violence gain to Stage~2; placing Stage~1 before it changes AP by a further \tMinusBzTMean{} points. During Stage~2, $s^{(1)}$ on a fixed probe batch did not change and no parameter outside $\mathrm{mlp}_1$ was updated.

\subsection{Robustness and Efficiency}
\label{sec:robust}

Under VadCLIP's official training setting, XD-Violence AP improves from \offBase{}$\pm$\offBaseSd{} to \offOurs{}$\pm$\offOursSd{}, with a paired gain of \offGainMean{}$\pm$\offGainSd{} points and \offGainPos{} positive runs. The improvement therefore persists beyond our protocol. Absolute values differ between the two settings, so each comparison is made within one setting.

The method adds \pvrParams{} parameters, increases inference latency by \cxUcfLatRel--\cxXdLatRel\%, and adds at most \cxUcfMemDelta{}~MB peak memory, measured on an \cxGpu{} per window of 256 snippets. This overhead keeps the detector suitable for lightweight surveillance deployment.

% Conclusion.

\section{Conclusion}
\label{sec:conclusion}

We showed that, in CLIP-based weakly supervised video anomaly detection, where refinement is applied matters as much as what is applied. Temporal variation extracted before fixed-length aggregation improves the classification output, whereas the same variation extracted afterwards degrades it. Semantic adaptation restricted to the semantic-side adapter improves the semantic output and leaves the classification pathway unchanged. Both gains are consistent across matched seeds, the XD-Violence gain persists under VadCLIP's official training setting, and the overhead is \pvrParams{} parameters and a few percent latency. These properties make the method a practical drop-in refinement for efficient surveillance and safety-monitoring systems.

\balance
\bibliographystyle{IEEEtran}
\bibliography{../refs/references}
\end{document}
